\section{Conclusion}\label{sec:conclusion}
We implemented an end-to-end sentiment system using 49,781 labeled IMDB reviews
with fixed, approximately balanced training, validation and test partitions.
The preparation pipeline performs and motivates tokenization, stop-word removal
and contextual lemmatization, protects negation, and preserves natural-text
inputs for the transformer. Quality checks, stable IDs and verified exports
provide a common basis for all team experiments.

The first experimental question is answered by the classical results: TF-IDF
outperforms mean-pooled Word2Vec with both LR and LinearSVC. The strongest
classical combination is TF-IDF + LinearSVC, with 88.684\% accuracy and 0.8865
F1. For the second question, LoRA-adapted DistilBERT achieves 92.564\% accuracy
and 0.9267 F1, improving accuracy by 3.880 percentage points and reducing errors
by 970 on the same 25,000 test reviews. The validation search favors 512-token
head-plus-tail inputs; adapter rank increases alone are less beneficial in the
recorded search. This is an observed whole-pipeline advantage under unequal
tuning, rather than proof of an isolated architectural effect.

The error analysis answers the third question by identifying mixed opinions
as the most common main category in the 100 sampled deep-model errors, alongside
possible label noise, sarcasm and negation. Truncation is relevant to a small
number of inspected cases but cannot explain every long-review error. The
Streamlit demo makes differences visible on new inputs and exposes a shared
sarcasm failure, while warning about long inputs and uncalibrated SVM scores.

The practical outcome is a strong contextual model with a compact trainable
update and an effective sparse linear baseline. Future work should prioritize
a strictly validation-only and equal-budget comparison, multiple seeds,
preprocessing and truncation ablations, independent error annotation and
out-of-domain evaluation. These steps would test the stability and scope of
the current findings without overstating the evidence already collected.
